% =====================================================================
% SMOTE25 paper draft, Method part (SMOTE family): working draft v2 (Claude, 28 Oct 2026)
% ---------------------------------------------------------------------
% P asked (28 Oct) for a draft of the paper from SMOTE (2002) to the present,
% written only from what the team has researched, centred on our part (the
% methods of the SMOTE family), with space left for the team members working
% on Datasets and on Performance Metrics. v1 (21 Oct) is replaced: all 37
% SMOTE-family papers were re-read in full on 28 Oct, which corrected the
% decision points of 21 of the 36 variants and added the class category.
% Markup, delete before submission (same macros as the Data section draft):
%  * \draftnum{}    number NOT checked against the PDFs (orange).
%  * \claudedraft{} sentence that interprets or argues (blue). Claude wrote it so
%                   the draft reads end to end; the interpretation must be P's
%                   own, so rewrite or delete every one.
%  * \evidence{}    where the numbers come from (grey).
%  * \classev{}     class-setting evidence of one paper, with page references (grey).
%  * \pcite{}       bibliography key not yet in references.bib (keys from
%                   Data section/out/section7_key_map.csv; wang2012multiclass is new).
%  * \pnote{}       a decision or check for P (boxed).
%  * \teamslot{}{}  section reserved for a team member (boxed).
% Generated parts (do not edit by hand; rerun Method section/build/make_tex.py
% SMOTE25_Method_Map_1.html "Method section/manuscript/generated" from the repo root):
%  * generated/smote-family-table.tex, generated/smote-family-methods.tex
% Sources: SMOTE25_Method_Map_1.html (76 papers; fields smote_rel, smote_via,
% class_category, class_evidence_pdf, smote_explain, all 37 SMOTE-family
% records checked against the PDFs on 28 Oct 2026);
% notes/smote25-development-brief.md (SMOTE 2002 facts with JAIR page numbers).
% Conventions: British spelling, Title Case headings, caption above tables,
% booktabs with \midrule only under the header, no em-dash.
% When moved into sections/ for Overleaf, change the \input paths below.
% =====================================================================
\providecommand{\draftnum}[1]{\textcolor{orange}{#1}}
\providecommand{\claudedraft}[1]{\textcolor{blue}{#1}}
\providecommand{\evidence}[1]{\par\noindent{\footnotesize\textcolor{gray}{\textbf{Evidence:} #1}}\par}
\providecommand{\classev}[1]{\par{\footnotesize\textcolor{gray}{\textbf{Class setting:} #1}}\par}
\providecommand{\pcite}[1]{[\textsc{#1}]}
\providecommand{\pnote}[1]{\par\noindent\fbox{\parbox{0.95\linewidth}{\small\textbf{[P]} #1}}\par}
\providecommand{\teamslot}[2]{\par\noindent\fbox{\parbox{0.95\linewidth}{\small\textbf{[Team: #1]} #2}}\par}
\providecommand{\fm}{$\bullet$}

% =====================================================================
\section{Introduction}
\label{sec:introduction}

\textcite{chawla2002smote} introduced the Synthetic Minority Over-sampling
Technique (SMOTE) in this journal in 2002. Instead of replicating minority
examples, SMOTE creates new ones by interpolating between a minority example
and one of its nearest minority neighbours. Sixteen years later,
\textcite{fernandez2018smote} reviewed the progress of SMOTE-based learning at
its 15-year anniversary. This survey follows the line from 2002 to the present
through a corpus of 76 papers and asks, for the methods that build on SMOTE,
what exactly they change, and in which class setting (two classes, several
classes, or both) they are shown to work.

Our account rests on three elements. First, we write the SMOTE generation rule
as six decisions and record which of them each method changes
(Section~\ref{sec:smote-by-decision}). Second, we sort the methods by the class
setting in which their papers actually generate synthetic data
(Section~\ref{sec:class-setting}). Third, we give for every method an account of
how it builds on SMOTE, with references to the pages of its paper
(Section~\ref{sec:methods}). All 37 SMOTE-family papers in the corpus were read
in full for this purpose.

\pnote{Motivation (growth of the literature since 2018, relation to Fern\'andez
et al.\ 2018, societal impact) and the research questions are still open; the
development brief has material for them. Keep this introduction to what the
Method part shows until the team agrees on the framing.}

% =====================================================================
\section{The SMOTE Generation Rule}
\label{sec:smote-rule}

Let $S_{\min}\subset\mathbb{R}^d$ be the minority examples of a two-class
training set and $S_{\mathrm{maj}}$ the majority examples. For a seed
$\mathbf{x}_i\in S_{\min}$, let $N_k(\mathbf{x}_i)$ denote its $k$ nearest
neighbours within $S_{\min}$ under the Euclidean distance. SMOTE creates a
synthetic example by choosing a neighbour $\mathbf{x}_j\in N_k(\mathbf{x}_i)$ at
random and placing a point at a random position on the segment that joins them
\parencite[pp.~328--329]{chawla2002smote}:
\begin{equation}
  \mathbf{x}_{\mathrm{new}} = \mathbf{x}_i + \lambda\,(\mathbf{x}_j-\mathbf{x}_i),
  \qquad \lambda\sim U(0,1).
  \label{eq:smote}
\end{equation}
The amount of oversampling is a percentage $N$: at $N=200\%$ each minority
example produces two synthetic points, towards two of its $k=5$ neighbours, and
amounts below $100\%$ apply the rule to a random subset of the minority class.
In the experiments SMOTE is combined with random undersampling of the majority
class at levels from $10\%$ to $2000\%$ \parencite[pp.~331, 337]{chawla2002smote}.
The same paper extends the rule to mixed features (SMOTE-NC) and to nominal
features (SMOTE-N), the latter without evaluation
\parencite[pp.~348--351]{chawla2002smote}.

Read as a procedure, Eq.~\eqref{eq:smote} fixes six choices in the simplest way
(Table~\ref{tab:decision-points}): every minority example is a seed with equal
weight; every seed receives the same number of points; the partner is drawn
uniformly from the five nearest minority neighbours; the point is drawn uniformly
on the segment; nothing is removed afterwards; and the generator runs once,
before training. We call these the six \emph{decision points} of the rule. A
general SMOTE-family generator can then be written as
\begin{equation}
  i\sim w(\cdot),\quad g_i,\quad \mathbf{x}_j\sim q(\cdot\mid\mathbf{x}_i),\quad
  \mathbf{x}_{\mathrm{new}}\sim r(\cdot\mid\mathbf{x}_i,\mathbf{x}_j),\quad F(\cdot),
  \label{eq:template}
\end{equation}
with a seed distribution $w$, an allocation $g_i$ of points to seeds (or to
clusters or classes), a partner distribution $q$, a placement distribution $r$,
and a filter $F$ applied to the enlarged training set; the sixth decision is
whether the procedure is called once or repeatedly inside a learner. SMOTE is the
special case with uniform $w$, constant $g_i=N/100$, uniform $q$ over
$N_5(\mathbf{x}_i)$, $r$ uniform on the segment, and the identity as $F$.

\begin{table}[t]
\centering\small
\caption{The Six Decision Points of the SMOTE Generation Rule. Papers: SMOTE-family
papers in the corpus that change each decision (37 papers; a paper is counted once
for every decision it changes).}
\label{tab:decision-points}
\setlength{\tabcolsep}{4pt}
\begin{tabular}{@{}p{0.16\linewidth}p{0.3\linewidth}p{0.24\linewidth}r@{}}
\toprule
Decision point & What it decides & SMOTE (2002) & Papers\\
\midrule
Seed selection & Which minority points become seeds, or how they are weighted & All seeds, equal weight & 24\\
Sample allocation & How many points each seed, cluster or class receives & $N/100$ per seed & 21\\
Neighbour choice & Which partner is paired with the seed & Uniform over 5 minority neighbours, Euclidean & 22\\
Placement rule & Where the new point lands & Uniform on the segment & 13\\
Post-filtering & What is removed or rejected after generation & Nothing & 12\\
Inside a pipeline & Whether generation runs inside boosting, bagging, a wrapper or training & Once, before training & 8\\
\bottomrule
\end{tabular}
\end{table}

\pnote{The published pseudo-code draws the gap once per attribute, not once per
point (p.~329), so a point lies in the box spanned by $\mathbf{x}_i$ and
$\mathbf{x}_j$ rather than on the segment. GDHS does the same on purpose (its
Eq.~3). Check the printed page; if confirmed, this is worth a footnote.}

% =====================================================================
\section{Review Method}
\label{sec:review-method}

\subsection{Corpus}
\label{sec:corpus}
The corpus is the team's collection of 76 papers, published from 1972 to 2026
and kept in four folders by problem setting: binary (24 papers), multiclass (21),
both (29) and surveys (2). Sixty-six papers propose a new method, five are
foundational (for example ENN, Tomek links and AdaBoost), three are empirical
studies and two are surveys. Each paper was recorded in a coding sheet (the
Method Map) with its approach, mechanism, baselines, datasets, classifiers,
metrics, stated limitations and code availability.

\teamslot{Review protocol}{Search strategy, databases, time window, inclusion and
exclusion criteria, and the PRISMA-style flow belong here (Section 3 of the
outline). The numbers above describe the corpus as it stands on 28 October 2026.}

\subsection{What Counts as a SMOTE-Family Method}
\label{sec:smote-definition}
We call a method \emph{SMOTE-family} when it creates synthetic minority examples
by interpolating between a minority seed and a neighbour as in
Eq.~\eqref{eq:smote}, or by a rule that extends it: a different law for
$\lambda$, a region built around the segment, extrapolation along the line
through the two points, or a convex combination of several instances. The
definition concerns the generator, not the name of a method. Each paper received
one of four labels: \emph{Core} (34 papers), where interpolation is the main
generator; \emph{Partial} (3: INOS, EMDO and GLOS), where it is one of several
generators; \emph{Study} (3), which evaluates or reviews SMOTE without proposing
a generator \pcite{fernandez2013analysing}\pcite{pradipta2021smote}\pcite{hassanat2022stop};
and \emph{Not counted} (12), related to SMOTE but generating in another way or not
at all. The remaining 24 papers have no direct link to SMOTE.

Ten of the twelve \emph{Not counted} papers generate synthetic minority examples,
but not between two instances: on the Mahalanobis contour of a seed (MDO, AMDO),
inside a sphere or ball around one point (MC-CCR, SHSampler, MC-NCLWO, MLOS),
along a random direction at a Gaussian distance (GDO), by perturbing points near a
cluster centroid (CBHSID), or by hill climbing on a potential field (RBO, MC-RBO).
The other two derive from SMOTE but drop its generator: RUSBoost replaces SMOTE in
SMOTEBoost by random undersampling, and AdaBoost.NC with random oversampling
\pcite{wang2012multiclass} duplicates examples. Two inclusions are judgement calls:
G-SMOTE samples in a hypersphere whose radius is set by the chosen neighbour and
which can be shrunk back to the segment, and GAMO replaces two-point interpolation
by a learned convex combination of all instances of a class. Both are counted
because they generalise the segment rather than replace it. Nine of the ten
region- or distribution-based generators appeared in 2018 or later.

\subsection{Coding Procedure}
\label{sec:coding}
All 37 SMOTE-family papers were read in full on 28 October 2026, including
supplementary material where the main paper left the class setting open. For each
paper we recorded (i) which of the six decision points the generator changes,
(ii) a short account of how it builds on SMOTE, and (iii) its class category,
defined by what the paper shows its generator doing:
\begin{itemize}
\item \emph{Binary only}: every experiment in which the generator creates
      synthetic data is a two-class problem. A multiclass dataset reduced to two
      classes (one class against the rest, merged classes, two classes kept and the
      rest dropped, or a ready-made split such as KEEL's \texttt{yeast-2\_vs\_4})
      counts here, and so does a multiclass extension that is described but never
      run.
\item \emph{Multiclass only}: the generator creates synthetic data with three or
      more classes kept, and is never run on a two-class problem.
\item \emph{Binary + Multiclass}: the same generator is run on at least one
      two-class problem and on at least one problem with three or more classes
      kept, without first reducing it to two classes.
\end{itemize}
A paper that generates more than one type of data (for example images and
tables) is classified for each data type separately
(Section~\ref{sec:other-data}), and its overall category combines them.
Experiments reported only in supplementary material count, since they show what
the method can do.
Every statement in the per-method accounts carries a page, section, table or
equation reference to its paper. The reading corrected the decision points of 21
of the 36 variants recorded in an earlier pass, mostly by adding decisions that
the first coding had missed.

\pnote{The reading was done by AI readers, one or two papers each, under fixed
rules (Method section/pdf-check/coding\_rules.md; their reports are the JSON files
in the same folder). State this in the paper as an assisted extraction and plan a
human check of a sample, as the Data section does.}

% =====================================================================
\section{Datasets}
\label{sec:datasets}

\teamslot{Datasets}{Which datasets the SMOTE-family papers use, where they come
from, how they are modified, and what the reader should check. The Data section
draft (Data section/manuscript/data-section-draft.tex) covers the whole corpus.
Two findings from the Method side belong here or should be cross-referenced:
(1) 23 of the 24 Binary-only SMOTE-family papers evaluate on multiclass sources
reduced to two classes, 15 by their own reduction and 8 through ready-made
two-class splits (mostly KEEL) whose reduction they do not describe
(Section~\ref{sec:class-setting}); (2) the same dataset name can hide different
class splits, so a multiclass name in a dataset list does not show multiclass
generation.}

% =====================================================================
\section{Performance Metrics}
\label{sec:metrics}

\teamslot{Performance metrics}{Which metrics the SMOTE-family papers report
(AUC, G-mean, F-measure, MCC, macro-averaged metrics for multiclass problems),
which statistical tests they run, and how comparable the results are. The
Section~7 draft (SMOTE25\_Evaluation\_Section.html) is the starting point. From
the Method side: several multiclass-capable methods score multiclass data with
macro-averaged or MAUC metrics (RSMOTE, MKC-SMOTE, GDHS), while Loyalty-SMOTE
scores its multiclass datasets with binary metrics.}

% =====================================================================
\section{The SMOTE Family by Decision Point}
\label{sec:smote-by-decision}

Table~\ref{tab:smote-family} lists the 37 SMOTE-family methods with their class
category and the decisions each one changes. A method is marked under every
decision it changes, so no method is forced into a single category. Variants
differ widely in how much they change: of the 36 variants, 8 change one decision,
8 change two, 8 change three, 9 change four, 2 change five and one (OREM with its
multiclass and boosting versions) changes all six.

\begin{table}[p]
\centering\footnotesize
\caption{SMOTE-Family Methods in the Corpus. Class: B binary only, M multiclass
only, B+M binary and multiclass (Section~\ref{sec:coding}; * provisional). Link:
Core or Partial. Decisions: Seed selection, sample Allocation, Neighbour choice,
Placement rule, post-Filtering, inside a Pipeline. Checked against the PDFs on
28 October 2026.}
\label{tab:smote-family}
\setlength{\tabcolsep}{3pt}
%% Generated by make_tex.py from SMOTE25_Method_Map_1.html (28 Oct 2026). Do not edit by hand.
\begin{tabular}{@{}p{0.33\linewidth}cclcccccc@{}}
\toprule
Method & Year & Class & Link & Seed & Alloc. & Nbr. & Place & Filter & Pipe\\
\midrule
SMOTE~\parencite{chawla2002smote} & 2002 & B & Core &  &  &  &  &  &  \\
SMOTEBoost~\pcite{chawla2003smoteboost} & 2003 & B+M & Core & \fm &  &  &  &  & \fm \\
SMOTE-Tomek, SMOTE-ENN~\pcite{batista2004study} & 2004 & B & Core &  &  &  &  & \fm &  \\
Borderline-SMOTE~\pcite{han2005borderline} & 2005 & B & Core & \fm &  & \fm & \fm &  &  \\
Cluster-SMOTE~\pcite{cieslak2006cluster} & 2006 & B & Core &  &  & \fm &  &  &  \\
ADASYN~\pcite{he2008adasyn} & 2008 & B & Core &  & \fm &  &  &  &  \\
LS~\pcite{cieslak2008ls} & 2008 & B & Core &  & \fm &  &  &  & \fm \\
Safe-Level-SMOTE~\pcite{bunkhumpornpat2009safelevel} & 2009 & B & Core & \fm &  &  & \fm &  &  \\
RAMOBoost~\pcite{chen2010ramoboost} & 2010 & B & Core & \fm & \fm &  &  &  & \fm \\
SSRBF, DSRBF~\pcite{fernandeznavarro2011dsrbf} & 2011 & B+M & Core &  & \fm &  &  &  & \fm \\
INOS~\pcite{cao2013inos} & 2013 & B & Partial &  & \fm &  &  & \fm &  \\
MWMOTE~\pcite{barua2014mwmote} & 2014 & B & Core & \fm & \fm & \fm &  &  &  \\
RBBoost~\pcite{diezpastor2015randombalance} & 2015 & B & Core & \fm & \fm &  &  &  & \fm \\
SCUT~\pcite{agrawal2015scut} & 2015 & M & Core &  & \fm &  &  &  &  \\
SMOTE-IPF~\pcite{saez2015smoteipf} & 2015 & B & Core &  &  &  &  & \fm &  \\
A-SUWO~\pcite{nekooeimehr2016asuwo} & 2016 & B & Core & \fm & \fm & \fm &  &  &  \\
CGMOS~\pcite{zhang2016cgmos} & 2016 & B & Core & \fm &  &  &  &  &  \\
SMOM~\pcite{zhu2017smom} & 2017 & M & Core &  &  & \fm &  &  &  \\
k-means SMOTE~\pcite{douzas2018kmeans} & 2018 & B & Core & \fm & \fm & \fm &  &  &  \\
G-SMOTE~\pcite{douzas2019gsmote} & 2019 & B & Core &  &  & \fm & \fm &  &  \\
GAMO~\pcite{mullick2019gamo} & 2019 & B+M & Core & \fm & \fm & \fm & \fm &  & \fm \\
AB-SMOTE~\pcite{almajzoub2020absmote} & 2020 & B & Core & \fm &  & \fm &  &  &  \\
EMDO~\pcite{yao2021emdo} & 2021 & M & Partial & \fm & \fm & \fm & \fm &  &  \\
HDE~\pcite{chen2021hde} & 2021 & B & Core &  &  & \fm & \fm &  & \fm \\
RSMOTE~\pcite{chen2021rsmote} & 2021 & B+M & Core & \fm & \fm &  &  &  &  \\
FW-SMOTE~\pcite{maldonado2022fwsmote} & 2022 & B & Core &  &  & \fm &  &  &  \\
KNNOR~\pcite{islam2022knnor} & 2022 & B & Core & \fm &  & \fm & \fm & \fm &  \\
SVDD-WSMOTE~\pcite{tao2022svddw} & 2022 & B & Core & \fm & \fm & \fm &  & \fm &  \\
GLOS~\pcite{han2023glos} & 2023 & M* & Partial & \fm & \fm & \fm & \fm &  &  \\
NROMM~\pcite{liu2023nromm} & 2023 & B+M & Core & \fm & \fm & \fm & \fm & \fm &  \\
OREM, OREM-M~\pcite{zhu2023orem} & 2023 & B+M & Core & \fm & \fm & \fm & \fm & \fm & \fm \\
MC-MBRC~\pcite{ma2024mcmbrc} & 2024 & B+M & Core & \fm & \fm & \fm &  & \fm &  \\
MKC-SMOTE~\pcite{wang2024mkc} & 2024 & M & Core & \fm &  & \fm & \fm & \fm &  \\
GDHS~\pcite{yan2025gdhs} & 2025 & B+M & Core & \fm & \fm & \fm &  & \fm &  \\
CISO~\pcite{yadav2025ciso} & 2026 & B & Core & \fm & \fm & \fm & \fm &  &  \\
Loyalty-SMOTE~\pcite{hu2026loyalty} & 2026 & B* & Core & \fm &  & \fm &  & \fm &  \\
QC-SMOTE~\pcite{upman2026qcsmote} & 2026 & B & Core & \fm & \fm &  & \fm & \fm &  \\
\addlinespace
Total (37) & & & & 24 & 21 & 22 & 13 & 12 & 8 \\
\bottomrule
\end{tabular}

\end{table}

\paragraph{Seed selection (24 papers).}
The most frequently changed decision is the choice of seeds, and the papers do
it in four ways. Eight concentrate generation near the class border or on hard
examples: Borderline-SMOTE~\pcite{han2005borderline} and
AB-SMOTE~\pcite{almajzoub2020absmote} use only ``danger'' seeds,
RAMOBoost~\pcite{chen2010ramoboost} draws seeds from a ranking that favours border
points, MWMOTE~\pcite{barua2014mwmote} and A-SUWO~\pcite{nekooeimehr2016asuwo}
weight seeds by closeness to the majority class, EMDO~\pcite{yao2021emdo} and
GLOS~\pcite{han2023glos} use borderline seeds for their interpolating generator,
and Loyalty-SMOTE~\pcite{hu2026loyalty} keeps boundary samples with a Loyalty
score between 0 and $m$. Seven favour safe or dense seeds:
k-means SMOTE~\pcite{douzas2018kmeans} keeps minority-dominated clusters,
KNNOR~\pcite{islam2022knnor} the densest minority points,
MC-MBRC~\pcite{ma2024mcmbrc} safe samples by fuzzy membership,
MKC-SMOTE~\pcite{wang2024mkc} points with at least half same-class neighbours,
GDHS~\pcite{yan2025gdhs} a safe factor, and CISO~\pcite{yadav2025ciso} and
QC-SMOTE~\pcite{upman2026qcsmote} safety or trust scores. Five remove noise or
type their seeds before generation: Safe-Level-SMOTE~\pcite{bunkhumpornpat2009safelevel}
skips pairs of noisy points, CGMOS~\pcite{zhang2016cgmos} weights seeds by the
expected gain in posterior certainty, RSMOTE~\pcite{chen2021rsmote} drops noise and
separates borderline from safe seeds, SVDD-WSMOTE~\pcite{tao2022svddw} removes
overlapped seeds inside the majority boundary, and NROMM~\pcite{liu2023nromm}
splits seeds into inland, borderline and trapped. In the last four the learner
sets the seeds: SMOTEBoost~\pcite{chawla2003smoteboost} and OREMBoost sample them
from the boosting distribution, RBBoost~\pcite{diezpastor2015randombalance} may
grow the majority class when it draws inverted class proportions, and
GAMO~\pcite{mullick2019gamo} has no single seed at all.

\paragraph{Sample allocation (21).}
Ten methods give seeds unequal numbers of points: ADASYN~\pcite{he2008adasyn} by
the share of majority neighbours, RAMOBoost, MWMOTE, A-SUWO, RSMOTE, SVDD-WSMOTE,
CISO and QC-SMOTE through weighted seed draws or quotas, k-means SMOTE by the
inverse density of clusters, and NROMM by local density. Seven set quotas per
class or region: SSRBF and DSRBF~\pcite{fernandeznavarro2011dsrbf} enlarge one
class at a time, SCUT~\pcite{agrawal2015scut} and GDHS fill classes to the mean
class size, MC-MBRC to the largest class minus its overlap, OREM-M~\pcite{zhu2023orem}
draws classes in proportion to their remaining quota, GLOS uses class
discreteness and EMDO splits each class's quota across ellipsoids. LS~\pcite{cieslak2008ls}
searches the SMOTE level per data segment with a wrapper, and RBBoost draws random
class proportions. INOS~\pcite{cao2013inos} splits the quota between an
interpolating and a Gaussian generator, and in GAMO the allocation is learned.

\paragraph{Neighbour choice (22).}
Nine methods restrict the partner to a region: the same cluster
(Cluster-SMOTE~\pcite{cieslak2006cluster}, MWMOTE, k-means SMOTE, NROMM), a
sub-cluster that does not cross the majority class (A-SUWO), other borderline
points (AB-SMOTE), the same ellipsoid under its Mahalanobis metric (EMDO), safe
and border samples of the same class (MC-MBRC) or the boundary set
(Loyalty-SMOTE). Five allow a majority point as partner or radius point, which the
2002 paper listed as future work \parencite[p.~348]{chawla2002smote}:
Borderline-SMOTE2, G-SMOTE~\pcite{douzas2019gsmote}, HDE~\pcite{chen2021hde},
OREM and GLOS. Four reweight the partners or directions (SMOM~\pcite{zhu2017smom},
GDHS, CISO, SVDD-WSMOTE), three replace the single partner by several instances
(KNNOR chains through all $k$ neighbours, MKC-SMOTE uses their centroid, GAMO all
instances of the class), and FW-SMOTE~\pcite{maldonado2022fwsmote} changes the
distance to a feature-weighted Minkowski metric.

\paragraph{Placement rule (13).}
Five methods change the law of $\lambda$ or the step: Borderline-SMOTE2
($\lambda\le 0.5$ towards majority partners), Safe-Level-SMOTE (shifted towards
the safer endpoint), OREM (half the step towards majority partners), CISO (a
truncated $q$-Gaussian on a segment between two neighbours of the anchor) and HDE
(a step from $N(0,1)$, which can extrapolate). Five leave the segment: G-SMOTE (a
truncated, deformed hypersphere), MKC-SMOTE (the line from the neighbours'
centroid through the seed and beyond), KNNOR (chained steps), GLOS (a step capped
at the nearest-neighbour distance, plus two non-interpolating generators) and EMDO
(points along ellipsoid axes in pure ellipsoids). GAMO uses a learned convex
combination, QC-SMOTE keeps the best of several candidates, and NROMM uses an
optimised embedding for trapped seeds and a constrained rule for categorical
values.

\paragraph{Post-filtering (12).}
Four methods clean the enlarged training set as a whole: SMOTE-Tomek and
SMOTE-ENN~\pcite{batista2004study}, SMOTE-IPF~\pcite{saez2015smoteipf}, INOS (noise
links) and Loyalty-SMOTE. Five test each candidate and reject or replace it:
KNNOR, SVDD-WSMOTE, QC-SMOTE, MKC-SMOTE (which over-generates and keeps the
candidates nearest the centroid) and OREM-M (which drops points nearest to another
minority class). Three remove or move majority examples in the regions where
synthetic points were placed: NROMM, MC-MBRC and GDHS.

\paragraph{Inside a pipeline (8).}
Eight methods call the generator repeatedly inside a learner: in every round of
AdaBoost.M2 (SMOTEBoost, RAMOBoost, RBBoost, OREMBoost), in every member of a
bagging ensemble (HDE), inside a wrapper that tunes the sampling levels (LS),
during the evolutionary training of a radial basis function network (DSRBF), or
inside adversarial training (GAMO).

\evidence{Method Map smote\_via after the 28 Oct PDF check (old values kept in
smote\_via\_v1); smote\_explain of each method. The groupings within each
decision are Claude's reading of those accounts.}

\pnote{The groupings in these paragraphs (for example border vs safe seeds) are
a new reading of the accounts, not fields of the Method Map. Keep only the ones
you want to defend, and check them against Section~\ref{sec:methods}.}

% Flush Table~\ref{tab:smote-family} (a full float page) here; otherwise every
% later table queues behind it until the end of the document. Drafting layout only.
\clearpage

% =====================================================================
\section{Binary, Multiclass, or Both}
\label{sec:class-setting}

\begin{table}[t]
\centering\small
\caption{SMOTE-Family Methods by Class Category and Period (37 Papers).}
\label{tab:class-period}
\begin{tabular}{lrrrr}
\toprule
Period & Binary only & Multiclass only & Binary + Multiclass & Total\\
\midrule
2002--2009 & 7 & 0 & 1 & 8\\
2010--2017 & 7 & 2 & 1 & 10\\
2018--2026 & 10 & 3 & 6 & 19\\
\addlinespace
All & 24 & 5 & 8 & 37\\
\bottomrule
\end{tabular}
\end{table}

Table~\ref{tab:class-period} shows that most SMOTE-family methods generate data
for two classes only. Twenty-four of the 37 are Binary only. Twenty-three of
these evaluate on multiclass sources reduced to two classes: 15 reduce the classes
themselves (one class against the rest, merged classes, or two classes kept and
the rest dropped) and 8 use ready-made two-class splits, mostly from KEEL, without
describing the reduction; the remaining one, Cluster-SMOTE, collapses the attack
types of its own network data into intrusive versus not. Eleven Binary-only papers
mention multiclass use without running it: four describe an extension (SMOTE,
ADASYN, k-means SMOTE, QC-SMOTE) and seven name it as future work (RAMOBoost,
MWMOTE, RBBoost, A-SUWO, KNNOR, SVDD-WSMOTE, CISO).

Methods that handle several classes are more recent. One of the 8 SMOTE-family
methods of 2002 to 2009 is Multiclass only or Binary + Multiclass (SMOTEBoost),
against 3 of 10 in 2010 to 2017 and 9 of 19 in 2018 to 2026. Five methods are
Multiclass only (SCUT, SMOM, EMDO, GLOS, MKC-SMOTE); their rules would cover two
classes as a special case, but their papers never show it. Eight are Binary +
Multiclass. In three of these the two-class evidence is limited: GAMO's two-class
runs are in its supplement, MC-MBRC's two-class data are synthetic, and OREM is
run on two-class data while its multiclass version OREM-M is run on multiclass
data.

\begin{table}[t]
\centering\small
\caption{How the 13 Multiclass-Capable SMOTE-Family Methods Handle Several Classes.}
\label{tab:mc-strategies}
\begin{tabular}{p{0.3\linewidth}p{0.6\linewidth}}
\toprule
Strategy & Methods\\
\midrule
To the largest class & SMOM, RSMOTE, MKC-SMOTE\\
To the mean class size & SCUT, GDHS\\
Class by class, against the larger classes & SSRBF/DSRBF, NROMM, OREM-M, MC-MBRC\\
Per-class generator or region & GAMO, EMDO, GLOS\\
Through the ensemble & SMOTEBoost (AdaBoost.M2 carries the classes)\\
\bottomrule
\end{tabular}
\end{table}

Table~\ref{tab:mc-strategies} groups the 13 methods that generate for several
classes by how they do it. None runs a separate binary generator for each
one-vs-one or one-vs-rest subproblem; several instead set each class against the
rest or against the larger classes within one procedure (EMDO, NROMM, OREM-M,
MC-MBRC), and RSMOTE uses one-vs-one only for its classifiers. Decomposition with binary SMOTE nevertheless remains a reference
point: in the study of \pcite{fernandez2013analysing}, one-vs-one with SMOTE was
among the best combinations.

\claudedraft{A dataset list is therefore a poor guide to what a method can do:
most papers that list multiclass datasets have reduced them to two classes, and a
multiclass extension stated in a paper is rarely tested.}

\pnote{Two categories are provisional. Loyalty-SMOTE runs on five multiclass
datasets but oversamples only the smallest class and scores with binary metrics
(coded Binary only, rule 1). GLOS's only two-class run is a parameter-tuning demo
on Sonar (coded Multiclass only). Decide both before the numbers above are final.}

\evidence{Method Map class\_category, class\_evidence\_pdf, class\_reduced and
mc\_strategy (28 Oct 2026); reduction split from class\_reduced, hand-coded in
make\_summary2.py (SELF 15, READY 8).}

\subsection{Papers on Other Data Types}
\label{sec:other-data}

Three of the 37 SMOTE-family papers generate data other than tables: INOS for
time series, GAMO for images, and KNNOR for tables and images. We classify each
data type on its own with the three rules of Section~\ref{sec:coding}
(Table~\ref{tab:other-data}). INOS is Binary only: its 48 tasks are built from
seven UCR datasets, four of them with 14 to 50 classes, by taking one class
against the rest \pcite{cao2013inos}. GAMO is Multiclass only on images, where all
of its seven datasets keep 5 to 50 classes, and Binary + Multiclass on tables, but
its tabular runs appear only in the supplement \pcite{mullick2019gamo}; its
overall category therefore rests on the supplement. KNNOR is Binary only on both:
its image experiments set male against female faces in UTKFace and airplanes
against the other nine CIFAR-10 classes, and show generated images without
classification results \pcite{islam2022knnor}.

\begin{table}[!ht]
\centering\small
\caption{SMOTE-Family Papers on Other Data Types, Classified per Data Type.}
\label{tab:other-data}
\begin{tabular}{@{}llp{0.18\linewidth}p{0.36\linewidth}l@{}}
\toprule
Method & Data type & Category for this data & Evidence & Overall\\
\midrule
INOS (2013) & Time series & Binary only & 48 two-class tasks from 7 UCR datasets, one class vs the rest (Sec.~4.2, pp.~7--8) & B\\
GAMO (2019) & Image & Multiclass only & MNIST, Fashion-MNIST, CIFAR10, SVHN (10 classes), CelebA, LSUN (5), SUN397 (50) (Secs.~4.1--4.4) & B+M\\
 & Tabular & Binary + Multiclass & Supplement only: Abalone19, IJCNN1, Magic (2 classes); Chess (18), Cover Type (7) (Supp.\ Tables~7--8) & \\
KNNOR (2022) & Tabular & Binary only & 19 two-class datasets, 18 from KEEL (Table~2, p.~11) & B\\
 & Image & Binary only & UTKFace male vs female; CIFAR-10 airplane vs the rest; images shown, no scores (Sec.~6, p.~14) & \\
\bottomrule
\end{tabular}
\end{table}

\claudedraft{Outside tables the SMOTE family is thin: three papers, of which only
GAMO generates for several classes, and only on images does it do so in its main
text.} Outside the family, the corpus has five more papers on other data types
(COG-OS and GDO on text, and Buda et al., DLexiBoost and NBU on images); they are
not classified here.

% =====================================================================
\section{How Each Method Builds on SMOTE}
\label{sec:methods}

This section gives, for each of the 37 methods, how its generator relates to
Eq.~\eqref{eq:smote}, the decisions it changes, and the evidence for its class
category. Methods are grouped by class category and ordered by year.

%% Generated by make_tex.py from SMOTE25_Method_Map_1.html (28 Oct 2026). Do not edit by hand.
%% Each account was written from a full reading of the PDF on 28 Oct 2026 (AI reader); check against the paper.
\subsection{Binary only (24 Methods)}
\label{sec:methods-B}
\begin{description}
\item[SMOTE (2002)] \parencite{chawla2002smote} This paper defines the reference rule and changes none of the six decision points: every minority sample is a seed with N/100 points each (a random N\% of seeds when N $<$ 100), the partner is drawn from the k = 5 minority nearest neighbours, and $\lambda \sim$ U(0,1) on the segment (Sec 4.2, p.328; algorithm, p.329). Generation runs once before training, combined with random undersampling of the majority (Sec 4.3, p.331). SMOTE-NC changes only the distance used to find neighbours and sets nominal values by a vote of the neighbours (Sec 6.1, p.348).
  \textit{Steps changed:} none (the reference rule).
  \classev{9 datasets, all used as two classes (Sec 5.1, pp.333-334; Table 2, p.334): Pima, Phoneme, Adult, E-state, Satimage, Forest Cover, Oil, Mammography, Can; Adult also run with SMOTE-NC (Sec 6.1, pp.348-349). The paper focuses on two-class problems; a multiclass use is stated but not run (pp.333-334). \textit{Reduction:} Satimage (6 classes): smallest class vs the other 5 merged (Sec 5.1 item 5, p.333). Forest Cover (7 classes): only Ponderosa Pine and Cottonwood/Willow kept, the other classes dropped (Sec 5.1 item 6, pp.333-334).}
\item[SMOTE-Tomek, SMOTE-ENN (2004)] \pcite{batista2004study} SMOTE first balances the training set (Sec. 5, p.25); then Tomek-link pairs are removed from both classes, or, in the ENN variant, every example misclassified by its 3 nearest neighbours is removed from both classes (Sec. 4.3, pp.23-24; Fig. 2). Seeds, neighbours and placement are as in SMOTE; the only addition is cleaning after generation.
  \textit{Steps changed:} Post-filtering.
  \classev{15 two-class datasets built from 13 UCI sources (Sec. 5, Table 3, pp.24-25); the paper considers only two-class problems (Sec. 4.1, p.22). C4.5, 10-fold CV, AUC (Tables 4-7). \textit{Reduction:} Smallest class vs the rest merged (Sec. 5, p.24), e.g. Vehicle van, New-thyroid hypo, E.Coli iMU, Satimage 4, Flag white, Glass, Nursery; Letter gives Letter-a (a vs rest) and Letter-vowel (all vowels vs rest); Splice gives two one-vs-rest sets (Table 3).}
\item[Borderline-SMOTE (2005)] \pcite{han2005borderline} For each minority point its m nearest neighbours in the whole training set are found (Sec. 3, Steps 1-2, p.882); points with only majority neighbours are treated as noise, and points with at least half majority neighbours form the DANGER set, which are the only seeds. In v1 each DANGER seed interpolates towards s of its k = 5 minority neighbours with $\lambda \sim$ U(0,1) as in SMOTE (Step 4, p.882). v2 also interpolates towards the seed's nearest majority neighbour with $\lambda$ in (0, 0.5), so those points stay near the seed (pp.882-883).
  \textit{Steps changed:} Seed selection, Neighbour choice, Placement rule.
  \classev{4 datasets used as two-class problems (Sec. 4, Table 2, p.883): Circle (simulated), Pima, Satimage (class 4 vs the rest), Haberman. The paper studies only two-class problems (Sec. 4, p.883). C4.5, 3 x 10-fold CV, TP rate and F-value (p.884). \textit{Reduction:} Satimage: class 4 vs all remaining classes merged (Sec. 4, p.883; Table 2).}
\item[Cluster-SMOTE (2006)] \pcite{cieslak2006cluster} k-means is run on the minority examples only and SMOTE is applied separately inside each cluster, after which original and synthetic minority points are merged (Sec. III-C, p.4). The interpolation partner therefore always comes from the seed's own cluster; seeds, the SMOTE rate per cluster and uniform placement on the segment are as in SMOTE (Sec. III-B, p.3; Sec. V-A, p.5).
  \textit{Steps changed:} Neighbour choice.
  \classev{Two datasets built by the authors from Notre Dame campus traffic (Sec. II, p.2): Packets (346,620 instances, 2,106 alerts) and Destinations (4,080 instances, 147 alerts) (Table I, p.2), both treated as intrusive vs not intrusive (Sec. II, p.2). Cluster-SMOTE was run only on Packets with 3 and 5 k-means clusters (Sec. V-A, Fig. 4, p.5); on Destinations the minority formed one cluster (Sec. V-B, pp.5-6). \textit{Reduction:} No benchmark multiclass dataset; SNORT attack-type labels are collapsed to intrusive vs not (Sec. II, p.2).}
\item[ADASYN (2008)] \pcite{he2008adasyn} ADASYN keeps the SMOTE generator: a partner is drawn from the seed's K nearest minority neighbours and the point is seed + $\lambda$(partner - seed), with K = 5 in the experiments (Eq. 5, p.1324; Sec. III-C, p.1325). What changes is allocation: the total G = (ml - ms) beta is shared out as $g_{i}$ proportional to the share of majority points among each seed's K neighbours (Eqs. 2-4, pp.1323-1324), instead of an equal number per seed.
  \textit{Steps changed:} Sample allocation.
  \classev{Algorithm defined for labels y in \{1,-1\} (Sec. II, p.1323). Artificial 2-class Gaussian set (Fig. 1, p.1324) and 5 UCI datasets used as two-class problems (Table I, Sec. III-A, pp.1324-1325): Vehicle, Pima, Vowel, Ionosphere, Abalone. A multiclass extension is described (Sec. III-D, p.1327) but never run. \textit{Reduction:} Vehicle (4 classes): van vs the rest (p.1324). Vowel (11 classes): first vowel vs the rest (p.1325). Abalone (29 classes): only class 18 vs class 9 kept (p.1325). Pima and Ionosphere are natively two-class.}
\item[LS (2008)] \pcite{cieslak2008ls} LS leaves the SMOTE generator unchanged but applies it per segment: a Hellinger-distance tree of height at most 2 splits the training data into up to 4 segments (Alg 1, Sec 3, pp4-5), and a 5-fold CV wrapper searches each segment for its own undersampling level (0-100\%) and SMOTE level (0-500\%) (Alg 2, Sec 4, p5). The number of synthetic points per seed therefore varies by segment, and generation runs inside a wrapper search rather than once before training. Seeds, partner and placement inside a segment are as in SMOTE.
  \textit{Steps changed:} Sample allocation, Inside a pipeline.
  \classev{20 datasets (Table 2, p7): real-world from Chawla et al. 2002, LIBSVM, the KDD'99 Intrusion subset and biological sets, plus UCI (Sec 5.3, pp7-8). All used as two-class problems with one minority percentage each (Table 2); Intrusion keeps u2r vs normal (Sec 5.3, p8). Motivating experiments use artificial two-class radial data (Sec 2, Fig 2-3, Table 1, pp2-4); Pima is the worked example (Fig 4, p6). \textit{Reduction:} Multiclass UCI sources reduced to smallest class vs the rest; Letter uses vowels vs consonants (Sec 5.3, p8). Which datasets were converted is not stated.}
\item[Safe-Level-SMOTE (2009)] \pcite{bunkhumpornpat2009safelevel} Each minority point has a safe level, the number of minority points among its k nearest neighbours (Eq 7, p478), and the ratio of the seed's and partner's safe levels bounds the gap (Fig 1, Sec 3, pp478-479): ratio 1 gives U(0,1) as in SMOTE, above 1 pulls the point towards the seed, below 1 towards the partner. When the partner's safe level is 0 the seed is copied, and when both are 0 nothing is generated, so noise pairs are skipped. The partner is still drawn uniformly from the 5 nearest minority neighbours and nothing is removed afterwards.
  \textit{Steps changed:} Seed selection, Placement rule.
  \classev{Studies only the two-class case (Sec 1, p475). Two UCI datasets (Table 3, p480): Satimage (6,435; 626 positive) and Haberman (306; 81 positive). k=5, 10-fold CV, C4.5, Naive Bayes and SVM (Sec 4, p480). \textit{Reduction:} Satimage: class 4 vs all other classes merged (Sec 4, p480).}
\item[RAMOBoost (2010)] \pcite{chen2010ramoboost} The generator is SMOTE's segment interpolation towards one of the seed's k2 nearest minority neighbours (Eq. 1, p.1626; Algorithm 1 step 5). In each boosting round, seeds are drawn with replacement from a ranking that weights each minority point by a logistic function of the number of majority points among its k1 neighbours (steps 2-4, Eqs. 2-3), so border points receive more synthetic samples. Generation is repeated in every AdaBoost.M2 round on a minority subset resampled by the boosting weights (steps 1, 6-10).
  \textit{Steps changed:} Seed selection, Sample allocation, Inside a pipeline.
  \classev{19 UCI/ELENA datasets, all run as two-class problems (Tables I-II, pp.1630-1631); only two-class problems considered (Sec. IV, p.1630); multiclass is future work (Sec. V, p.1641). \textit{Reduction:} One-vs-rest (e.g. Wine 1 vs 2+3, Vehicle van, Segment brickface, Letter Z vs A-Y), merged classes (Texture, Page\_blocks) or two classes kept and the rest dropped (Abalone 18 vs 9, Yeast POX vs CYT) (Table II, p.1631).}
\item[INOS (2013)] \pcite{cao2013inos} INOS splits the synthetic quota between two generators: 70\% from ESPO, a regularised Gaussian of the minority covariance that does not interpolate (Sec. 3.2.2, Eq. 13; r = 0.7, Sec. 4.3), and 30\% from ADASYN in a signal space with the null space removed (Sec. 3.3, Eqs. 14-15). The ADASYN part allocates points per seed by the share of negatives among its neighbours and interpolates uniformly towards one of the 5 nearest positives, as in SMOTE. Afterwards, Tomek-like noise links between synthetic positives and nearest negatives are removed (Sec. 3.4, pp.5-6).
  \textit{Steps changed:} Sample allocation, Post-filtering.
  \classev{Restricted to two-class time series (p.2; Sec. 3, p.3). 48 two-class apportions from 7 UCR datasets (Adiac, FaceAll, 50words, SLeaf, Wafer, Yoga, TwoPats), each a separate binary task (Sec. 4.2, pp.7-8, Table 1). \textit{Reduction:} Adiac (37 classes), FaceAll (14), 50words (50), SLeaf (15): one class positive, the rest merged as negative, one apportion per chosen class (Sec. 4.2, p.8).}
\item[MWMOTE (2014)] \pcite{barua2014mwmote} Seeds are limited to an informative minority set near the borderline majority points, after dropping minority points whose neighbours are all majority (Alg. 1 steps 1-6; Sec. 4.1, p.5). Each seed gets a weight from its closeness to the majority class and its density, and seeds are drawn in proportion to it, so counts per seed vary (Alg. 1 steps 7-9, 12a; Sec. 4.2, pp.6-7). The partner is a random member of the seed's minority cluster from average-linkage clustering, not one of its 5 nearest neighbours (Alg. 1 steps 10, 12b; Secs. 4.3-4.4). Placement is uniform on the segment and nothing is removed afterwards.
  \textit{Steps changed:} Seed selection, Sample allocation, Neighbour choice.
  \classev{4 artificial two-class domains (Sec. 5.3, pp.11-12) and 20 UCI datasets used as two-class problems (Sec. 5.4, p.13, Table 5); networks with 2 output neurons (Sec. 5.2). Multiclass use named as future work (Sec. 6, p.19). \textit{Reduction:} Multiclass sources converted to two-class problems with the chosen minority and majority classes listed in Table 5 (Sec. 5.4, p.13); the table could not be extracted from the text.}
\item[RBBoost (2015)] \pcite{diezpastor2015randombalance} Each ensemble member gets a training set of the original size with a random class split (Alg. 1 lines 6-7, p.98): the class that must grow is filled with SMOTE (k = 5; Sec. 3, p.97) and the other is randomly undersampled, and when the split is inverted SMOTE grows the original majority class (Alg. 1 lines 12-15; Sec. 4.1, p.98). Partner and placement are as in SMOTE. Generation is redone for every member of a plain or bagging ensemble, or in every AdaBoost.M2 round (RB-Boost, Alg. 2, Sec. 4.2, p.100).
  \textit{Steps changed:} Seed selection, Sample allocation, Inside a pipeline.
  \classev{Stated for two-class problems (Abstract; Sec. 1, p.97); Algorithms 1-2 use labels y in \{-1,+1\} (pp.98, 100). Two-class 2-D Gaussian simulation (Sec. 5, p.100). 86 binary datasets: 20 HDDT (Table 2, p.102) and 66 KEEL (Table 3, p.103), 5x2-fold CV, J48 (Sec. 6, p.101). Multiclass extension listed as future work only (Sec. 7, p.110). \textit{Reduction:} 66 KEEL sets derived from 16 sources (footnote 7, p.101) as one-vs-rest or merged-group splits (e.g. ecoli-0-1-4-7\_vs\_2-3-5-6, yeast-2\_vs\_4, page-blocks-1-3\_vs\_4); HDDT sets such as letter, pendigits, satimage are presumably binarised, reduction not described.}
\item[SMOTE-IPF (2015)] \pcite{saez2015smoteipf} Generation is plain SMOTE with k = 5, oversampling to a 50/50 balance (Sec. 3.2, p.189). The only change is a filter applied afterwards: the Iterative-Partitioning Filter, an ensemble of C4.5 models on 9 partitions with majority voting, repeated until fewer than 1\% of examples are removed in 3 rounds, which removes noisy original and synthetic examples of both classes (Sec. 3.1, pp.188-189; Sec. 3.3, p.190).
  \textit{Steps changed:} Post-filtering.
  \classev{30 synthetic two-class 2-D datasets (Sec. 4.1, pp.190-191) and 9 real datasets, each one minority class vs the rest, 8 from KEEL (Sec. 4.2, p.191; Table 1, p.192), with class or attribute noise added. \textit{Reduction:} Multiclass sources turned into two classes by joining one or more classes as positive and the rest as negative (Sec. 4.2, p.191), likely ecoli, newthyroid and cleveland.}
\item[A-SUWO (2016)] \pcite{nekooeimehr2016asuwo} Minority points are grouped by a complete-linkage clustering that will not merge two minority sub-clusters across a majority sub-cluster (Sec. 3.1, pp.26-28). Each sub-cluster's quota is proportional to its cross-validated classification error (Sec. 3.2, Eqs. 3.2-3.3), seeds are drawn with weights favouring closeness to the majority class (Sec. 3.3, Eqs. 3.4-3.9), and the partner is one of the seed's nearest neighbours in the same sub-cluster (Alg. 3.1, p.33). Placement stays uniform on the segment (Eq. 3.10).
  \textit{Steps changed:} Seed selection, Sample allocation, Neighbour choice.
  \classev{Drive PDF is the 2016 dissertation, Chapter 3 being the A-SUWO paper. Defined for binary problems (p.24; Sec. 3.5, p.46). 16 datasets, each a minority class vs the rest (Table 3.1, p.35); multiclass use only future work (Sec. 3.5, p.46). \textit{Reduction:} 12 multiclass sources reduced to named class(es) vs all other classes (e.g. Vehicle van, Ecoli pp, Glass 1, Segment WINDOW) (Table 3.1, p.35).}
\item[CGMOS (2016)] \pcite{zhang2016cgmos} CGMOS keeps SMOTE's interpolation step unchanged (Sec. 2.3, p.1625). Seeds are drawn in proportion to a weight: the average relative change in kernel-density posterior certainty over all training points if a minority point were added at that seed (Eqs. 5-6, p.1625). Sec. 3 proves that the expected certainty gain is at least that of SMOTE on the training data (Theorem 1).
  \textit{Steps changed:} Seed selection.
  \classev{Set up for two classes only (Sec. 2, p.1624; Sec. 5, p.1629). 30 UCI datasets used as two-class problems (Sec. 4.1, p.1626; Table 1, p.1627). \textit{Reduction:} Multiclass datasets keep the smallest class as minority and merge the rest (Sec. 4.1, p.1626); which datasets is not stated.}
\item[k-means SMOTE (2018)] \pcite{douzas2018kmeans} k-means clusters the whole input space, and only clusters whose majority-to-minority ratio is below a threshold are kept, so seeds are minority points inside minority-dominated clusters (Sec. 3.1, PDF pp.7-8; Algorithm 1). Samples are shared across kept clusters in proportion to their sparsity, so sparse clusters get more (Sec. 3.1, PDF pp.8-9). Inside each cluster plain SMOTE runs on that cluster's points only, so partners come from the same cluster; placement is unchanged and there is no filtering. With one cluster and no threshold the method reduces to SMOTE (Sec. 3.2, PDF p.9).
  \textit{Steps changed:} Seed selection, Sample allocation, Neighbour choice.
  \classev{All experiments two-class; a multiclass extension is said to be possible but not covered (Sec. 1, PDF p.4). 90 datasets (Abstract): 12 UCI sets plus randomly undersampled copies, 19 KEEL two-class sets and 10 simulated MADELON sets (Sec. 4.4, PDF p.14; Table 2, PDF p.15), plus 5 two-class 2-D toy sets (Figs. 10-14). \textit{Reduction:} Multiclass UCI sources (e.g. breast tissue, ecoli, glass, iris, libra, segment, vehicle, wine) turned into one-vs-rest problems with the smallest class as minority (Sec. 4.4, PDF p.14); KEEL sets are ready-made binary splits (Table 2).}
\item[G-SMOTE (2019)] \pcite{douzas2019gsmote} Seeds and allocation are SMOTE-like (Sec. 4.1, PDF p.9). The surface point that sets the radius is chosen by a selection strategy: a random minority neighbour as in SMOTE, the nearest majority point, or whichever of the two is closer (Algorithm 1, PDF p.8; Sec. 4.3 step 4). The segment is replaced by a point drawn in a hypersphere around the seed with that radius, truncated and deformed into a hyper-spheroid along the seed-to-surface axis (Sec. 4.2; Sec. 4.3 steps 6-9, PDF pp.14-16). With full truncation and deformation and the minority strategy the region becomes SMOTE's segment (Sec. 5.4, PDF p.24).
  \textit{Steps changed:} Neighbour choice, Placement rule.
  \classev{All experiments two-class: 28 base datasets (13 UCI, 13 KEEL, 2 simulated MADELON) plus undersampled copies, 69 in total (Sec. 5.1, PDF p.18; Table 1, pp.18-19); binary AUC, F-score and G-mean (Sec. 5.2). \textit{Reduction:} Multiclass sources (e.g. iris, wine, glass, vehicle, segmentation, ecoli, libras, vowel, led, dermatology) appear as two-class sets in Table 1; how they were reduced is not stated.}
\item[AB-SMOTE (2020)] \pcite{almajzoub2020absmote} As in Borderline-SMOTE, only danger minority points (at least half but not all of their M nearest neighbours are majority) are seeds (Sec. II.B, p.32; Sec. III, pp.32-33). The difference is that the k nearest neighbours, and so the partner, are searched only within the borderline set (Sec. III, Figs. 5-6, p.33). Placement is P'i + gap (P'j - P'i) with gap in [0,1], as in SMOTE (Sec. III, p.33).
  \textit{Steps changed:} Seed selection, Neighbour choice.
  \classev{Framed for binary classification (title; Sec. I, p.31). Eight two-class datasets (Table I, p.33): Abalone9-18, Flare, Wine Quality (KEEL); Haberman (UCI); Customer Churn Wireless Telecom (CrowdAnalytix); Telco Customer Churn, Employee Attrition, Sales Win/Loss (IBM). J48, 5-fold CV (Sec. IV, pp.33-34). \textit{Reduction:} No reduction described; Abalone9-18, Flare and Wine Quality are KEEL two-class versions of multiclass sources.}
\item[HDE (2021)] \pcite{chen2021hde} The partner is one of the k = 5 nearest majority instances (nearest enemies), found with a mix of Euclidean distance and how often earlier ensemble members disagree on the two points (Sec. 3.3, Eqs. 3-4, pp.6-7). The step is x + r($x_{i}$ - x) with r drawn from N(0,1), so points cluster near the seed and can fall on either side of it (Eq. 5, p.8). Generation runs inside every member of a sequential bagging ensemble (Secs. 3.1-3.4).
  \textit{Steps changed:} Neighbour choice, Placement rule, Inside a pipeline.
  \classev{Targets binary tasks (Abstract; Sec. 2; Conclusion). 42 binary sets: 28 KEEL and 14 LibSVM (Sec. 4.1, Table 1, pp.9-10). \textit{Reduction:} KEEL x\_vs\_y and one-vs-rest splits of multiclass sources (Table 1); reduction not described.}
\item[FW-SMOTE (2022)] \pcite{maldonado2022fwsmote} The only change is how the k nearest minority neighbours are found: Euclidean distance is replaced by a feature-weighted Minkowski distance whose weights come from a supervised feature ranking, and low-ranked features can be dropped from the distance (Eqs. 4-5, 9-12, pp.4-5; Algorithm 1 steps 2-9). Seeds, the number per seed, the random partner and placement on the segment are as in SMOTE (Algorithm 1 steps 10-14), and the method reduces to SMOTE with equal weights and p = 2 (p.6).
  \textit{Steps changed:} Neighbour choice.
  \classev{Defined for labels \{-1,+1\} (Sec. 3.3, p.5). 42 two-class datasets: 30 low-dimensional UCI/KEEL and 12 high-dimensional microarray sets (Table 1, p.7; Sec. 4.1, pp.6-7); binary AUC and G-mean. \textit{Reduction:} Low-dimensional sets are mostly KEEL 'x vs y' and one-vs-rest versions of multiclass sources (Table 1, p.7); the paper does not describe the reduction.}
\item[KNNOR (2022)] \pcite{islam2022knnor} Seeds are only the densest d\% of minority points, ranked by distance to their k-th minority neighbour, with d chosen automatically (Alg. 1, Sec. 3.1, pp.3-4). Instead of one random partner, the point moves step by step through all k nearest minority neighbours, each step bounded by the smallest majority-minority distance (Eq. 5, Alg. 2, Sec. 3.2, pp.6-8). A candidate is kept only if all its k nearest neighbours are minority points (Alg. 3, Sec. 3.3, p.8).
  \textit{Steps changed:} Seed selection, Neighbour choice, Placement rule, Post-filtering.
  \classev{19 two-class datasets (18 KEEL plus CM1; Table 2, Sec. 4.1, p.11) and two-class image demos (UTKFace; CIFAR-10 airplane vs rest; Sec. 6, p.14). Multiclass use only future work (Sec. 7, p.16). \textit{Reduction:} KEEL splits of multiclass UCI sources (Sec. 4.1, p.11); CIFAR-10 one-vs-rest.}
\item[SVDD-WSMOTE (2022)] \pcite{tao2022svddw} Noisy or overlapped minority points inside the majority SVDD sphere are removed before generation (Secs. 4.1-4.2). Each remaining seed gets a number of points proportional to a weight combining closeness to the boundary and sparsity (Eqs. 21-24), and the partner is drawn from its 5 nearest neighbours in proportion to their weights (Eq. 25, p.25). Placement is SMOTE's; synthetic points that fall inside the SVDD boundary are rejected and regenerated, and any still missing are filled by copying the seed (Alg. steps 36-52, pp.26-27).
  \textit{Steps changed:} Seed selection, Sample allocation, Neighbour choice, Post-filtering.
  \classev{Four 2-D simulated two-class sets (Sec. 5.1), 24 binary datasets from 16 UCI sources (Table 2, pp.34-35) and bearing-fault data with fault types against normal (Table 8, pp.44-45). Defined for one majority and one minority class (Eqs. 7-8, 24); multiclass is future work (Sec. 6). \textit{Reduction:} One-vs-others or class-vs-class (Abalone 16 vs 6, Yeast3-5) (Sec. 5.3; Table 2).}
\item[CISO (2026)] \pcite{yadav2025ciso} Anchors are drawn with replacement in proportion to a safety score, so each anchor gets a varying number of points (Eqs. 8-9, p.8). Two different minority neighbours of the anchor are drawn with weights from their clearance and a q-Gaussian locality kernel (Eqs. 10-11). The new point lies on the segment between those two neighbours, not between anchor and neighbour, with $\lambda$ drawn from a truncated q-Gaussian instead of U(0,1) (Eqs. 12-14, p.12). Nothing is removed afterwards. Running time is 39 times SMOTE's at the median (Sec. 6.7, p.29).
  \textit{Steps changed:} Seed selection, Sample allocation, Neighbour choice, Placement rule.
  \classev{Binary labels only (Sec. 3.1, p.5). 37 binary KEEL datasets for evaluation, 5 KEEL datasets for tuning and 8 larger UCI datasets (Sec. 5.1, p.15; Tables 1-2). Multiclass is future work (Sec. 7, p.34). \textit{Reduction:} Multiclass sources made binary as smallest class vs rest (Sec. 5.1; Table 1 caption), e.g. Waveform, Statlog Landsat, Page Blocks, Letter; KEEL sets are ready-made binary versions.}
\item[Loyalty-SMOTE (2026)] \pcite{hu2026loyalty} Generation is plain SMOTE (k = 5) applied only to boundary minority samples whose Loyalty, an inverse-distance-weighted balance of same-class and other-class neighbours, lies between 0 and m (Sec. 3.3). Samples with negative Loyalty are treated as possible noise and kept out of generation, so the partner pool is also restricted. Afterwards Loyalty is recomputed: flagged minority samples still below 0 are removed and flagged majority samples are removed by an LOF test (Sec. 3.4, Eqs. 10-12).
  \textit{Steps changed:} Seed selection, Neighbour choice, Post-filtering.
  \classev{30 binary UCI/KEEL datasets (group 1, Table 2) plus other two-class groups, and 5 multiclass UCI datasets (group 4, Table 5: yeast 10 classes, HR 3, employee 3, Drug200 4, MEDV 9) with results in Sec. 4.4.4 and Tables 15-18. \textit{Reduction:} Group 1 includes KEEL x\_vs\_y splits and one-vs-rest versions of multiclass sources (Table 2); reduction not described. \textit{Several classes:} On multiclass data the largest class is the majority and the smallest the minority (Sec. 4.4.4); SMOTE oversamples only the smallest class (Sec. 3.3); low-Loyalty samples are relabelled to the class with the highest Attraction before generation (Eq. 9, Sec. 3.2). Whether the middle classes stay in the data is read differently by two checks; only one class is ever oversampled, and results use binary metrics.}
  \pnote{Provisional: the PDF check read this as Binary + Multiclass; coded Binary only under rule 1 (only the smallest class is ever oversampled, binary metrics). Decision for P. Category decision needed: only the smallest class is oversampled on multiclass data, scored with binary metrics.}
\item[QC-SMOTE (2026)] \pcite{upman2026qcsmote} Seeds are drawn in proportion to a trust score combining safe level, minority density and distance from the majority class (Eqs. 2-6, Sec. 3.2, p.5), and the number of attempts per seed also follows trust (Eq. 7, Sec. 3.3, p.6). Each candidate is still $x_{i}$ + $\lambda$($x_{j}$ - $x_{i}$) with a 5-NN minority partner, but $\lambda \sim$ U(0, $\alpha_{r}$) shrinks with overlap and the best of $K_{r}$ candidates is kept by a purity and clearance score (Eqs. 8-18, Secs. 3.4-3.5, pp.6-7). A candidate that fails the thresholds is replaced by a copy of the most trusted point near the seed (Eqs. 19-21, Sec. 3.6, p.7).
  \textit{Steps changed:} Seed selection, Sample allocation, Placement rule, Post-filtering.
  \classev{30 binary datasets from the VS-SMOTE benchmark (KEEL-style), IR 1.82 to 32.73 (Table 1, p.8; Sec. 4.1, p.8). CatBoost, repeated stratified CV, oversampling on training folds only (Secs. 4.2-4.3, p.8). A one-vs-rest use is described (Sec. 3.1, p.4; Sec. 3.7, p.7) but never run. \textit{Reduction:} Uses ready-made KEEL two-class splits (one-vs-rest such as ecoli1-4, glass0/1/6, vehicle0-3, yeast1/3/4/5; merged-class splits such as glass-0-1-2-3\_vs\_4-5-6, yeast-1-2-8-9\_vs\_7) (Table 1, p.8); the paper does not describe the reduction.}
\end{description}
\subsection{Multiclass only (5 Methods)}
\label{sec:methods-M}
\begin{description}
\item[SCUT (2015)] \pcite{agrawal2015scut} Synthetic points come from standard SMOTE run inside each under-represented class (Fig. 1); the only change is how many points each class gets, enough to reach the mean class size m with a rate computed per class (Sec. 2). Seeds, partners within the class and placement are as in SMOTE, and generation is a one-off step before training. The EM cluster-based undersampling acts on real points of the large classes, not on synthetic ones.
  \textit{Steps changed:} Sample allocation.
  \classev{8 multiclass datasets with all classes kept: 7 KEEL (Thyroid 3 classes, Lymphography 4, Pageblocks 5, Dermatology 6, Autos 6, Ecoli 8, Yeast 10) and UCI Wine Quality (7) (Sec. 3.1, Table 1). No two-class experiment. 10-fold CV with J48, SMO, Naive Bayes, IBk (Sec. 3). \textit{Several classes:} Mean class size m is the common target: classes below m are SMOTEd up to m with a per-class rate, classes above m are reduced to m by EM clustering and near-uniform sampling across clusters (Sec. 2, Fig. 1; Sec. 3.1). No decomposition, no ensemble.}
\item[SMOM (2017)] \pcite{zhu2017smom} SMOM keeps SMOTE-style interpolation between a seed and one of its k1 = 12 neighbours (Supp. Secs. 3, 5). For trapped instances, each neighbour direction gets a selection weight that penalises directions whose segment sphere contains instances of other classes (Algorithm S1; Supp. Sec. 3), so directions are drawn by weight rather than uniformly. NBDOS clustering decides which instances are outstanding (uniform directions) or trapped (weighted directions).
  \textit{Steps changed:} Neighbour choice.
  \classev{Read from the supplement only: the main PDF in Drive contains the INOS paper, not SMOM. 27 multiclass datasets, UCI and ordinal-regression benchmarks, all classes kept, with per-group minority and majority metrics (Supp. Tables S3-S10, S17-S20). No two-class experiment appears in the supplement. \textit{Several classes:} Each minority class is processed in turn against all other classes and oversampled until all classes have the same size (Supp. Sec. 7, Table S21); neighbours are searched over all classes; no decomposition or ensemble.}
  \pnote{Main PDF mislinked in Drive; placement law and allocation not confirmed from the main text.}
\item[EMDO (2021)] \pcite{yao2021emdo} Each minority class is covered by Mahalanobis ellipsoids and its quota is split across them by inverse density times purity (Eqs. 32-34, p.11). In pure ellipsoids points are drawn along the ellipsoid axes without interpolation (Eqs. 37-40). In impure ellipsoids a modified Borderline-SMOTE interpolates: only borderline points are seeds, and the partner comes from their 5 Mahalanobis neighbours within the ellipsoid (Eqs. 41-42, p.12).
  \textit{Steps changed:} Seed selection, Sample allocation, Neighbour choice, Placement rule.
  \classev{9 KEEL/UCI sets with all classes kept (3 to 6 classes; Table 1, p.13) and 2 sensor sets (Shuttle 7, Mafalda 3; Table 6, p.16). Only two-class data is a toy illustration (Fig. 3). In 5 of the 11 sets only one class is oversampled. \textit{Several classes:} Any class whose size ratio to the largest class is below the threshold is grown to N\_max/IR (Eq. 1, p.4); each minority class gets its own ellipsoids fitted against all other classes (Eqs. 15-18).}
\item[GLOS (2023)] \pcite{han2023glos} Only one of GLOS's three generators interpolates: AMBO, used for borderline seeds (Alg. 2, p.2076). Its partner is any instance inside a radius-R sphere, possibly of another class, and the step is capped at the seed's nearest-neighbour distance rather than covering the whole segment (Alg. 2 line 12). The number of points per seed comes from discreteness and is split across classes by their discreteness ratios. The other two generators place points on Mahalanobis contours (MMO) or at random within a radius (RRO).
  \textit{Steps changed:} Seed selection, Sample allocation, Neighbour choice, Placement rule.
  \classev{21 multiclass UCI/KEEL datasets with 3 to 11 classes kept (Sec. 4.1, Table 1). One parameter-tuning run on Sonar (Fig. 7, p.2079) also reports GLOS results; the paper does not state Sonar's class count (two classes in the standard UCI version). \textit{Several classes:} Classes are labelled minority or majority by comparing their discreteness with the dataset-wide value (Eqs. 1, 2, 5); minority seeds are typed safe, borderline or small-disjunct and each type gets its own generator (MMO, AMBO, RRO) (Algs. 1-3).}
  \pnote{Provisional: Multiclass only; would be Binary + Multiclass if the two-class Sonar tuning run (Fig. 7) counts. Decision for P.}
\item[MKC-SMOTE (2024)] \pcite{wang2024mkc} A minority point is a seed only if at least half of its 5 nearest neighbours, taken from all classes, share its class (Sec. III step 2; Alg. 1). Instead of a random minority partner, the method uses the centroid of the seed's neighbours and places new points on the line from that centroid through the seed, possibly beyond the seed (Sec. III steps 3-5; Alg. 1 lines 23-26). It generates twice the needed number and keeps the candidates nearest the centroid (Sec. III step 6, p.5).
  \textit{Steps changed:} Seed selection, Neighbour choice, Placement rule, Post-filtering.
  \classev{21 UCI datasets, all multiclass (Abstract; Sec. IV-A-1, p.5; Table 1), all classes kept and scored with multiclass MG and MAUC (Eqs. 6-7, p.6); 5-fold CV, CART and SVM. No two-class experiment. \textit{Several classes:} Every class with N\_max/N\_i $>$ 1.45 is oversampled towards the largest class size, each class independently, with the neighbour search over all classes (Alg. 1; Sec. III, pp.4-5). No decomposition.}
\end{description}
\subsection{Binary + Multiclass (8 Methods)}
\label{sec:methods-BM}
\begin{description}
\item[SMOTEBoost (2003)] \pcite{chawla2003smoteboost} The generator is unchanged SMOTE (Sec 2, p.109), but it runs inside every AdaBoost.M2 round and the synthetic points are discarded after the round (Sec 3, p.110; Fig. 1, p.111). New points are created by sampling seeds from the round's boosting distribution, which the authors say focuses generation on misclassified or boundary minority cases (pp.110, 116). Allocation, partner, placement and the absence of filtering are as in SMOTE.
  \textit{Steps changed:} Seed selection, Inside a pipeline.
  \classev{4 datasets (Sec 4, Table 2, p.112): KDDCup-99 with all 5 classes kept (DoS, Probe, Normal, U2R, R2L); Mammography and Phoneme (two-class); Satimage (reduced to two). Results tabulated only for U2R (Table 3, p.113); R2L omitted for space. RIPPER base learner (p.111). \textit{Reduction:} Satimage: smallest class vs the rest merged (p.112). KDDCup-99 not reduced; only DoS and Normal subsampled for size (p.112). \textit{Several classes:} AdaBoost.M2 core over labels \{1..C\} (Fig. 1, p.111); on KDDCup-99 SMOTE is applied separately to each minority class (U2R at 100/300/500\%, R2L at 100\%) (p.112); the other three classes are not oversampled.}
\item[SSRBF, DSRBF (2011)] \pcite{fernandeznavarro2011dsrbf} The generator is unchanged SMOTE: every original pattern of the chosen class is a seed joined to one of its k = 5 nearest same-class neighbours (Sec. 1, p.1821; Sec. 3.4, p.1826). What changes is which class is oversampled and how many times, set by class size (SSRBF) or by the lowest per-class sensitivity of the current best network (DSRBF). In DSRBF generation runs inside the training loop; SSRBF runs before training.
  \textit{Steps changed:} Sample allocation, Inside a pipeline.
  \classev{13 datasets (Sec. 4.1, p.1827; Table 3, p.1828): 5 two-class (Hepatitis, BreastC, Haberman, CYTvsPOX, ME2vsOther) and 8 with all classes kept (Newthyroid 3, Balance 3, Saureus4 4, Anneal 5, Glass 6, Zoo 7, E. coli 8, Yeast 10). SSRBF and DSRBF run on all 13 (Table 4) with one softmax RBF network, no decomposition. \textit{Reduction:} CYTvsPOX and ME2vsOther are derived from Yeast (Sec. 4.1, p.1827). \textit{Several classes:} One class is oversampled at a time: SSRBF enlarges the currently smallest class J times before training; DSRBF pauses the memetic training at set generations and oversamples the class with the lowest sensitivity, each time by its original size (Secs. 3.3-3.4, pp.1825-1827).}
\item[GAMO (2019)] \pcite{mullick2019gamo} GAMO creates each sample as a convex combination of all training points of a minority class (Eq. 1, Sec. 3.2, p.1697), not as a point on a seed-neighbour segment, so there is no single seed and no k-NN partner: a generator network learns the weights adversarially to place points near class peripheries that the classifier gets wrong. Generation runs inside end-to-end training that alternates generator, classifier and discriminator updates (Alg. 1, p.1698), usually in a learned feature space.
  \textit{Steps changed:} Seed selection, Sample allocation, Neighbour choice, Placement rule, Inside a pipeline.
  \classev{Main paper: seven image datasets, all multiclass with selected classes kept: MNIST and Fashion-MNIST (10 classes, IR 100; Sec. 4.1, Table 1), CIFAR10 and SVHN (10 classes; Sec. 4.2, Table 2), CelebA and LSUN (5 classes, small and large; Sec. 4.3, Table 3), SUN397 (50 classes; Sec. 4.4, Table 4). Supplement (Sec. 4, Tables 7-8, p.8): GAMO also run on three two-class sets (Abalone19, IJCNN1, Magic) and two multiclass sets with all classes kept (Chess 18 classes, Cover Type 7). \textit{Reduction:} None in the main paper. Supplement uses Abalone19, the KEEL two-class version of Abalone (reduction not described); IJCNN1 and Magic are two-class by nature. \textit{Several classes:} The c-1 minority classes are each filled up to the majority count (Sec. 3.1, p.1697); a shared class-conditional unit feeds one instance-generation unit per minority class, which outputs convex weights over that class's instances (Sec. 3.2, Eq. 1); labels of generated points are drawn in proportion to (Pc - Pi) (Alg. 1). One c-class classifier, no decomposition.}
\item[RSMOTE (2021)] \pcite{chen2021rsmote} Minority points whose k neighbours are all majority are dropped as noise, and the remaining seeds are split into borderline and safe groups by 2-means on their relative density (Defs. 1-3, pp.401-402; Algorithm 1, p.403). Each group gets a share matching its size, and within a group each seed gets a weight (k-m)/(k+1), so safer seeds produce more points (Step 3, p.402). Points are made with SMOTE's own routine, uniform on the segment, once before training (Algorithm 1 line 22).
  \textit{Steps changed:} Seed selection, Sample allocation.
  \classev{Two-class: 4 synthetic 2-D sets (Sec. 4.3.1, p.406) and 27 UCI sets used as binary (Table 2, p.409), each at 0-30\% label noise. Multiclass: 8 UCI sets with all classes kept, 3 to 5 classes (Table 7, p.416; Sec. 4.3.3, pp.413-414), macro-averaged metrics (Eqs. 7-10, p.405). Same generator in both. \textit{Reduction:} Several of the 27 binary sets come from multiclass UCI sources (e.g. ecoli, iris, wine, segment, vehicle; Table 2, p.409); the reduction is not described but counts suggest one class vs the rest. \textit{Several classes:} Every class smaller than the largest is oversampled with RSMOTE up to the largest class size (Sec. 4.3.3, p.413); classification uses diversified one-against-one classifiers (p.414).}
\item[NROMM (2023)] \pcite{liu2023nromm} Clustering splits minority samples into inland, borderline and trapped (Sec. 3.2). Inland and borderline seeds are interpolated, inland ones with a partner from the same cluster and borderline ones with a nearby majority neighbour; trapped seeds get points from an optimised embedding, not interpolation (Sec. 3.3, Eqs. 2-4). The number of points per seed follows (1 - local density) times the class-size gap, categorical values come from a VDM-constrained set near the seed (Sec. 4.1), and borderline and trapped majority samples are removed afterwards.
  \textit{Steps changed:} Seed selection, Sample allocation, Neighbour choice, Placement rule, Post-filtering.
  \classev{20 KEEL multiclass datasets with 3 to 10 classes kept (Sec. 5.1; Table 2, p.7) and 10 two-class KEEL datasets (D1-D10) in the validity check (Sec. 5.3, p.9; Table 6; Supp. Table 1). \textit{Reduction:} D1-D10 are KEEL two-class versions of multiclass sources (one-vs-rest such as Yeast4, Ecoli4; merged splits such as Ecoli-3-5 vs 0-6-7) (Supp. Table 1). \textit{Several classes:} One-versus-ensemble: classes sorted by size, each non-largest class oversampled against the original and synthetic data of all larger classes (Sec. 4.2.1; Fig. 4); afterwards inland samples of the overlap region are removed (Sec. 4.2.2).}
  \pnote{Core vs Partial debatable: trapped seeds (6-84\% per dataset) are not interpolated. The binary and multiclass variants appear to be run on separate data.}
\item[OREM, OREM-M (2023)] \pcite{zhu2023orem} OREM keeps every minority point as a seed with equal counts (Alg. 1). It replaces the k-NN partner set with 'assistant seeds' from a candidate region that grows until q = 5 consecutive majority points are met, keeping a candidate only if the sphere on the seed-candidate segment holds no majority point, so borderline majority points can be partners (Alg. 1; Sec. 3.1.2). The point is $x_{i}$ + g($x_{s}$ - $x_{i}$) with g $\sim$ U[0,1], halved when the partner is a majority point. OREM-M adds class-quota sampling and a nearest-neighbour rejection filter (Alg. 2), and OREMBoost runs the generator inside each AdaBoost.M2 round (Alg. 3).
  \textit{Steps changed:} Seed selection, Sample allocation, Neighbour choice, Placement rule, Post-filtering, Inside a pipeline.
  \classev{OREM on 28 two-class UCI datasets (Sec. 4.1.1, Table 2, p.6174); OREM-M on 21 multiclass datasets with 3 to 10 classes kept (Sec. 4.2.1, Table 6, p.6177); OREMBoost on both (Sec. 4.3, Tables 10-11). \textit{Reduction:} 20 of OREM's 28 sets come from multiclass sources reduced to one or merged classes vs the rest (Table 2). \textit{Several classes:} OREM-M (Alg. 2, Sec. 3.2, p.6172): only classes with imbalance degree above 1.2 are oversampled; one point at a time, a class is drawn in proportion to its remaining quota, and a synthetic point is rejected if its nearest neighbour belongs to another minority class. Native multiclass, not OvO/OvR.}
  \pnote{Variants: OREM is B, OREM-M is M, OREMBoost is BM; the paper is BM through OREMBoost and the union.}
\item[MC-MBRC (2024)] \pcite{ma2024mcmbrc} Seeds are only the safe samples of each class, after membership-based noise removal (Alg. 1; Sec. 3). The partner is one of the 5 same-class neighbours among the safe and border samples, and placement is SMOTE's $x_{i}$ + U(0,1)($x_{j}$ - $x_{i}$) (Alg. 2). Each seed gets a random number of points up to the class quota divided by the number of safe seeds. Afterwards, a CCR-style energy sphere moves overlapping and generated samples of larger classes towards their class centre (Algs. 3-4, p.6).
  \textit{Steps changed:} Seed selection, Sample allocation, Neighbour choice, Post-filtering.
  \classev{20 multiclass UCI/KEEL datasets with 3 to 10 classes kept (Sec. 4.2, Table 2, p.7); 60 synthetic two-class datasets (Sec. 5.2, Fig. 5) and 50 more in the ablation (Sec. 5.4, Fig. 7). \textit{Several classes:} Membership to each class centre splits every class into safe, overlapping and noisy samples (Alg. 1); every class gets a quota N\_max minus its overlap count (Alg. 2); cleaning goes through classes by descending size, treating all larger classes as majority (Alg. 4). No decomposition.}
\item[GDHS (2025)] \pcite{yan2025gdhs} Generation interpolates between a minority seed and one of its 5 same-class neighbours (Eq. 3, p.3), but each seed-direction pair gets a weight that favours safe seeds (same-class share among neighbours, Eq. 2) and directions whose synthetic region holds few other-class points (Eqs. 5-10). Each minority class is filled to the mean class size by weighted random draws rather than a fixed number per seed (Alg. 1). Afterwards, original majority points in the synthetic regions, or overlapping other majority classes, are removed (Sec. 3.2, Eqs. 11-16).
  \textit{Steps changed:} Seed selection, Sample allocation, Neighbour choice, Post-filtering.
  \classev{20 datasets (Sec. 4.1, Table 1, pp.6-7): 14 KEEL multiclass sets with all classes kept (3 to 10 classes, e.g. car, glass, yeast, zoo) and 6 two-class software-defect sets (camel-1.4, camel-1.6, hive-0.12.0, velocity-1.5, xalan-2.4, xerces-1.2). All three variants run on all 20 sets (Tables 5-8, pp.10-11). \textit{Reduction:} None; the defect sets are natively two-class. \textit{Several classes:} Mean class size m is the target (Eq. 1): classes below m are oversampled up to m, classes above m are cleaned by one of three strategies (Sec. 3.2, pp.5-6). No decomposition.}
\end{description}


% =====================================================================
\section{What the Method View Shows}
\label{sec:method-findings}

\emph{Changes accumulate.} The mean number of decisions changed per variant rose
from 1.7 in 2002 to 2009 (7 variants) to 2.0 in 2010 to 2017 (10) and 3.6 in 2018
to 2026 (19). Twelve of the 36 variants change four or more decisions.
\claudedraft{When a method changes the seeds, the partner, the placement and the
filter together, a comparison with SMOTE alone cannot tell which change produced
the gain.}

\emph{The partner, the placement and the filter changed most after 2018.} Since
2018, 17 of 19 variants change the neighbour choice (5 of 17 before), 11 change
the placement rule (2 of 17 before) and 9 add a filter after generation (3 of 17
before). Generation inside a pipeline did not grow (2, 3 and 3 variants in the
three periods).

\emph{The questions of 2002 have mostly been taken up, with one exception.} The
2002 paper named four directions for future work
\parencite[p.~348]{chawla2002smote}: other ways of creating synthetic points
(13 methods change the placement rule), a focus on misclassified examples (the
border-oriented seed rules and the boosting methods), and the use of majority
neighbours (five methods allow a majority partner). The fourth, an adaptive choice
of $k$, is not the subject of any method in the corpus; MKC-SMOTE still lists it as
future work.

\pnote{``Not the subject of any method'' rests on the Method Map fields; the PDF
check was not asked about $k$. Confirm before keeping it.}

\emph{Recent variants cost more than SMOTE.} CISO reports a median resampling
time of 47 ms per fold against 1.2 ms for SMOTE, about 39 times
\pcite{yadav2025ciso}, and QC-SMOTE about 39.47 times (its Table~7)
\pcite{upman2026qcsmote}. CISO also reports no significant improvement over SMOTE,
class weighting or random oversampling in PR-AUC.
\claudedraft{The move towards safer, more elaborate generators has not yet been
matched by evidence that the extra computation pays off;
Section~\ref{sec:metrics} examines that evidence.}

\emph{The family has a boundary, and it is moving.} Since 2018, nine new
generators in the corpus have left interpolation for spheres, contours, potential
fields or fitted distributions, and deep generative models (CTGAN, SOS, GOIO,
DGOT) learn the minority distribution instead of interpolating between examples.

\evidence{Method Map smote\_via by year (variants: 7, 10, 19 per period);
CISO Sec.~6.7, p.~29; QC-SMOTE Sec.~4.5.5, Table~7, p.~10; CISO's PR-AUC result from
its limitations field (\draftnum{not re-checked}).}

% =====================================================================
\section{Recommendations for Describing a New Variant}
\label{sec:method-recommendations}
\begin{enumerate}
\item State which of the six decisions the method changes, and keep the others as
      in SMOTE unless there is a reason not to.
\item When more than one decision changes, report a comparison that changes one
      decision at a time against SMOTE.
\item State the class setting in which the generator is run, and whether any
      multiclass dataset was reduced to two classes and how.
\item State how nominal and mixed features are handled, or that they are not.
\item Report the running time relative to SMOTE on the same data.
\end{enumerate}

\pnote{Keep, cut or reword; each item should point back to a finding above.
Item~3 overlaps with the Data section's recommendations, items~2 and~5 with
Section~7.}

% =====================================================================
\section{Discussion and Conclusion}
\label{sec:conclusion}

\teamslot{Discussion and conclusion}{To be written once the Datasets and
Performance Metrics sections are in: what the field has done since 2002, what
remains open, and what the three parts (methods, data, evaluation) say together.}
